\documentclass[11pt,a4paper]{ctexart}
\usepackage[margin=2.0cm]{geometry}
\usepackage{amsmath,amssymb}
\usepackage{booktabs}
\usepackage{array}
\usepackage{hyperref}

\hypersetup{colorlinks=true,linkcolor=blue,urlcolor=blue}
\setlength{\parindent}{0pt}
\setlength{\parskip}{0.45em}

\title{代码中 Strang Splitting 离散公式整理}
\author{}
\date{}

\begin{document}
\maketitle

\section*{1. 总体说明}

代码中的深度方向 \(x\) 是 marching 变量。记一个完整深度步长为
\(\Delta x\)，代码变量为 \texttt{grid.dt}。主粒子和二次粒子都使用同一个
Strang splitting 框架：
\begin{equation}
  \boxed{
  \psi^{n+1}
  =
  \mathcal E_{\Delta x/2}
  \mathcal T_{\Delta x/2}
  \mathcal A_{\Delta x}
  \mathcal T_{\Delta x/2}
  \mathcal E_{\Delta x/2}
  \psi^n
  }.
\end{equation}

对应代码调用顺序：
\begin{equation}
  \text{energy half}
  \rightarrow
  \text{transport half}
  \rightarrow
  \text{angle full}
  \rightarrow
  \text{transport half}
  \rightarrow
  \text{energy half}.
\end{equation}

论文公式对应关系如下：
\begin{equation}
\begin{array}{lll}
\mathcal E & \text{energy/DG/CN step} & \text{Eq. (15)} \\
\mathcal T & \text{lateral }(y,z)\text{ transport step} & \text{Eq. (18), Eq. (22)} \\
\mathcal A & \text{angular diffusion step} & \text{Eq. (19), Eq. (20), Eq. (21)} \\
\text{primary/secondary split} & \psi=\psi_p+\psi_s & \text{Eq. (9), Eq. (10a), Eq. (10b)}
\end{array}
\end{equation}

当前代码中 CN 使用位置为：
\begin{equation}
\begin{array}{llll}
\text{子步} & \text{论文公式} & \text{是否使用 CN} & \text{代码实现方式}\\
\hline
\mathcal E & \text{Eq. (15)} & \text{是} & \text{DG/CN 半隐式能量递推}\\
\mathcal T & \text{Eq. (18), Eq. (22)} & \text{否} & \text{显式 MUSCL predictor--corrector}\\
\mathcal A & \text{Eq. (19)--(21)} & \text{是} & \text{sine/DST 对角化后的 CN 谱更新}
\end{array}
\end{equation}
注意：论文 Remark 3 说明 Eq. (22) 的 depth discretization 也可以采用 CN，并可显式处理 slope；
但当前代码没有按这个隐式/CN 形式实现 transport，而是使用显式二阶 predictor--corrector。

\section*{2. 代码变量与论文字母}

\begin{center}
\small
\begin{tabular}{lll}
\toprule
代码变量 & 论文/数学符号 & 含义 \\
\midrule
\texttt{F} & \(\psi^1_{i,j,g,m}\) & DG 能量第一个系数，或分布主系数 \\
\texttt{f\_F} & \(\psi^2_{i,j,g,m}\) & DG 能量第二个系数/斜率系数 \\
\texttt{F1}, \texttt{f\_F1} & \(\psi^1_{s},\psi^2_{s}\) & 二次粒子对应两个 DG 系数 \\
\texttt{dt} & \(\Delta x\) & 深度步长，transport/energy 半步传入 \(\Delta x/2\) \\
\texttt{dg} & \(\Delta E\) & 能量组宽度 \\
\texttt{dy}, \texttt{dz} & \(\Delta y,\Delta z\) & 横向网格距 \\
\texttt{du}, \texttt{dv} & \(\Delta u,\Delta v\) & 角变量网格距 \\
\texttt{Omend} & \((\Omega_{m,y},\Omega_{m,z})\) & 横向角方向 \\
\texttt{S\_s[g]} & \(S_g\) & stopping power \\
\texttt{T\_c[g]} & \(T_g\) & straggling 系数 \\
\texttt{sigma\_c[g]} & \(\sigma_{c,g}\) & catastrophic loss 截面 \\
\texttt{sig\_trg[g]} & \(\Sigma_{\mathrm{tr},g}\) & angular transport cross section \\
\texttt{density} & \(\rho\) & 材料密度因子 \\
\bottomrule
\end{tabular}
\end{center}

\section*{3. 主粒子 Strang Splitting}

主粒子对应论文 Eq. (9), Eq. (10a)。代码中每一步为
\begin{equation}
\begin{aligned}
  \psi_{p}^{(1)}
    &= \mathcal E_{\Delta x/2}(\psi_p^n),\\
  \psi_{p}^{(2)}
    &= \mathcal T_{\Delta x/2}(\psi_p^{(1)}),\\
  \psi_{p}^{(3)}
    &= \mathcal A_{\Delta x}(\psi_p^{(2)}),\\
  \psi_{p}^{(4)}
    &= \mathcal T_{\Delta x/2}(\psi_p^{(3)}),\\
  \psi_{p}^{n+1}
    &= \mathcal E_{\Delta x/2}(\psi_p^{(4)}).
\end{aligned}
\end{equation}
代码对 \(\psi^1\) 和 \(\psi^2\) 两个 DG 系数都执行同样的 transport 和 angle 子步。

\section*{4. 二次粒子 Strang Splitting}

二次粒子对应论文 Eq. (9), Eq. (10b)。代码结构同主粒子，但 energy step 使用由主粒子构造的
source \(Q_p\)：
\begin{equation}
\begin{aligned}
  \psi_{s}^{(1)}
    &= \mathcal E_{\Delta x/2}^{Q_p}(\psi_s^n;\psi_p),\\
  \psi_{s}^{(2)}
    &= \mathcal T_{\Delta x/2}(\psi_s^{(1)}),\\
  \psi_{s}^{(3)}
    &= \mathcal A_{\Delta x}(\psi_s^{(2)}),\\
  \psi_{s}^{(4)}
    &= \mathcal T_{\Delta x/2}(\psi_s^{(3)}),\\
  \psi_{s}^{n+1}
    &= \mathcal E_{\Delta x/2}^{Q_p}(\psi_s^{(4)};\psi_p).
\end{aligned}
\end{equation}

二次源离散求和在代码中为
\begin{equation}
  Q_{g,m}^{1}(i,j)
  =
  \Delta E
  \sum_{g'}\sum_{m'}
  \left[
    K^{E,1}_{g,g'} K^{\Omega}_{g',m',m}
    \max(\psi^{1}_{p,g',m'}(i,j),0)
    +
    \frac13
    K^{E,2}_{g,g'} K^{\Omega}_{g',m',m}
    \psi^{2}_{p,g',m'}(i,j)
  \right].
\end{equation}
这是代码对 Eq. (10b) 中 catastrophic secondary source 的离散实现。

\section*{5. Energy Step: Eq. (15)，使用 CN}

energy step 对每个固定的 \((i,j,m)\) 沿能量组 \(g\) 从高能到低能递推。
设当前旧值为
\begin{equation}
  P_g=\psi^{1,n}_g,\qquad R_g=\psi^{2,n}_g,
\end{equation}
新值为
\begin{equation}
  \widehat P_g=\psi^{1,n+1}_g,\qquad
  \widehat R_g=\psi^{2,n+1}_g.
\end{equation}
在代码里
\begin{equation}
  S_g=S_s[g],\quad S_{g+1}=S_s[g+1],\quad
  \sigma_g=\max(\sigma_{c,g},0),\quad
  h=\Delta x_{\mathcal E},\quad \Delta E=\texttt{dg}.
\end{equation}
其中 \(h=\Delta x/2\)。

代码使用的 Eq. (15) DG/CN 更新可写成下列代数形式。这里是当前实现中第一处使用
Crank--Nicolson 的地方；每个 Strang 步中调用两次，且每次传入半步 \(h=\Delta x/2\)。
先定义
\begin{equation}
  D_g
  =
  \frac1h
  +\frac12\rho\frac{S_g}{\Delta E}
  +\frac12\sigma_g,
  \qquad
  A_g
  =
  \frac{3}{2}\rho\frac{S_{g+1}}{\Delta E}\frac{1}{D_g},
\end{equation}
\begin{equation}
  B_g
  =
  \frac12\rho\frac{S_g}{\Delta E}\frac{1}{D_g},
  \qquad
  C_g
  =
  \frac12\rho\frac{S_{g+1}}{\Delta E}\frac{1}{D_g}.
\end{equation}
并记
\begin{equation}
  U_g=P_{g+1}-R_{g+1},\qquad V_g=P_g-R_g.
\end{equation}
若启用 straggling，代码加入
\begin{equation}
  \mathcal S_g
  =
  \frac12\rho\frac{T_{g+1}}{\Delta E^2}P_{g+1}
  +
  \frac12\rho\frac{T_{g-1}}{\Delta E^2}P_{g-1}
  -
  \rho\frac{T_g}{\Delta E^2}P_g.
\end{equation}
若未启用，则 \(\mathcal S_g=0\)。

定义
\begin{equation}
  M_g
  =
  \frac12\rho\frac{S_{g+1}}{\Delta E}U_g
  -
  \frac12\rho\frac{S_g}{\Delta E}V_g
  -
  \frac12\sigma_g P_g
  +
  \frac{P_g}{h}
  +
  \mathcal S_g
  +
  Q_g^1.
\end{equation}
则
\begin{equation}
  R^{(4)}_g=\frac{M_g}{D_g}.
\end{equation}
第二个 DG 系数满足
\begin{equation}
  \widehat R_g=\frac{N_g}{C_g^{R}},
\end{equation}
其中
\begin{equation}
\begin{aligned}
  N_g={}&
    R_g/h-\frac12\sigma_gR_g
    +\frac32\rho\frac{S_{g+1}}{\Delta E}U_g
    +\frac32\rho\frac{S_g}{\Delta E}V_g \\
  &-\frac32\rho\frac{S_g+S_{g+1}}{\Delta E}P_g
    -\frac12\rho\frac{S_{g+1}-S_g}{\Delta E}R_g
    -A_gM_g+Q_g^2 \\
  &+H_g(\widehat P_{g+1}-\widehat R_{g+1}),
\end{aligned}
\end{equation}
\begin{equation}
  C_g^{R}
  =
  \frac1h+\frac12\sigma_g
  +\frac32\rho\frac{S_g}{\Delta E}
  +\frac12\rho\frac{S_{g+1}-S_g}{\Delta E}
  +A_g\frac12\rho\frac{S_g}{\Delta E}.
\end{equation}
这里
\begin{equation}
  H_g
  =
  \frac32\rho\frac{S_{g+1}}{\Delta E}
  -
  A_g\frac12\rho\frac{S_{g+1}}{\Delta E}.
\end{equation}
最后
\begin{equation}
  \widehat P_g
  =
  B_g\widehat R_g
  +
  R^{(4)}_g
  +
  C_g(\widehat P_{g+1}-\widehat R_{g+1}).
\end{equation}
主粒子时 \(Q_g^1=Q_g^2=0\)。二次粒子时 \(Q_g^1,Q_g^2\) 来自上面的 secondary source。

\subsection*{5.1 Energy Step 的实际更新顺序}

上式不是同时求解 \(\widehat P_g\) 和 \(\widehat R_g\)，而是沿能量组从高能到低能反向扫掠。
原因是当前组 \(g\) 的新值依赖高能邻居的新值差
\begin{equation}
  \Delta_{g+1}
  =
  \widehat P_{g+1}-\widehat R_{g+1}.
\end{equation}
代码初始化最高能边界为
\begin{equation}
  \widehat P_{N_g}=0,
  \qquad
  \widehat R_{N_g}=0,
\end{equation}
对应 \texttt{next\_F = 0.0} 和 \texttt{next\_f = 0.0}。

对每个能量组 \(g=N_g-1,N_g-2,\ldots,0\)，代码执行：
\begin{equation}
  \Delta_{g+1}
  =
  \widehat P_{g+1}-\widehat R_{g+1},
\end{equation}
\begin{equation}
  \widehat R_g
  =
  \frac{
    N_g^{0}+H_g\Delta_{g+1}
  }{
    C_g^R
  },
\end{equation}
\begin{equation}
  \widehat P_g
  =
  B_g\widehat R_g
  +
  R_g^{(4)}
  +
  C_g\Delta_{g+1}.
\end{equation}
其中 \(N_g^0\) 是不含高能新值差 \(\Delta_{g+1}\) 的分子部分。

写成代码变量就是
\begin{equation}
\begin{array}{lll}
\widehat R_g &\leftrightarrow& \texttt{f\_val},\\
\widehat P_g &\leftrightarrow& \texttt{F\_val},\\
\Delta_{g+1} &\leftrightarrow& \texttt{hi\_delta = next\_F - next\_f},\\
B_g &\leftrightarrow& \texttt{F\_from\_f[idx]},\\
C_g &\leftrightarrow& \texttt{F\_from\_hi[idx]},\\
R_g^{(4)} &\leftrightarrow& \texttt{rhs4[idx]},\\
N_g^0 &\leftrightarrow& \texttt{rhs\_base[idx]},\\
H_g &\leftrightarrow& \texttt{rhs\_high\_coeff[idx]},\\
C_g^R &\leftrightarrow& \texttt{coe[idx]}.
\end{array}
\end{equation}

对应代码更新为
\begin{equation}
  \texttt{f\_val}
  =
  \frac{
    \texttt{rhs\_base[idx]}
    +
    \texttt{rhs\_high\_coeff[idx]}\cdot\texttt{hi\_delta}
  }{
    \texttt{coe[idx]}
  },
\end{equation}
\begin{equation}
  \texttt{F\_val}
  =
  \texttt{F\_from\_f[idx]}\cdot\texttt{f\_val}
  +
  \texttt{rhs4[idx]}
  +
  \texttt{F\_from\_hi[idx]}\cdot\texttt{hi\_delta}.
\end{equation}
随后代码令
\begin{equation}
  \texttt{next\_F}=\texttt{F\_val},
  \qquad
  \texttt{next\_f}=\texttt{f\_val},
\end{equation}
供下一个低能组 \(g-1\) 使用。

\subsection*{5.2 Energy Step 与三对角矩阵的关系}

如果只看能量方向的二阶离散，例如把能量 straggling 项
\(\partial_E(T\partial_E\psi)\) 完全隐式处理，则固定一个 \((i,j,m)\) 后，
未知量
\begin{equation}
  \widehat{\boldsymbol\psi}
  =
  \left[
  \widehat\psi_0,\widehat\psi_1,\ldots,\widehat\psi_{N_g-1}
  \right]^T
\end{equation}
会满足典型三对角线性系统
\begin{equation}
  \begin{bmatrix}
  d_0 & u_0 &        &        & 0\\
  \ell_1 & d_1 & u_1 &        &  \\
        & \ell_2 & d_2 & \ddots &  \\
        &        & \ddots & \ddots & u_{N_g-2}\\
  0     &        &        & \ell_{N_g-1} & d_{N_g-1}
  \end{bmatrix}
  \begin{bmatrix}
  \widehat\psi_0\\
  \widehat\psi_1\\
  \widehat\psi_2\\
  \vdots\\
  \widehat\psi_{N_g-1}
  \end{bmatrix}
  =
  \begin{bmatrix}
  b_0\\
  b_1\\
  b_2\\
  \vdots\\
  b_{N_g-1}
  \end{bmatrix}.
\end{equation}
这里
\begin{equation}
  \ell_g
  \quad\text{来自 }g-1\text{ 组耦合},\qquad
  d_g
  \quad\text{来自本组 }g,\qquad
  u_g
  \quad\text{来自 }g+1\text{ 组耦合}.
\end{equation}
这就是通常所说的 energy 方向三对角矩阵结构。

但是当前 CUDA 代码没有组装上面的三对角矩阵。当前实现中 straggling 项
\begin{equation}
  \mathcal S_g
  =
  \frac12\rho\frac{T_{g+1}}{\Delta E^2}P_{g+1}
  +
  \frac12\rho\frac{T_{g-1}}{\Delta E^2}P_{g-1}
  -
  \rho\frac{T_g}{\Delta E^2}P_g
\end{equation}
使用的是旧值 \(P_{g+1},P_g,P_{g-1}\)，因此它进入右端项，而不是进入新值矩阵。
新值部分只保留对高能邻居差值
\(\widehat P_{g+1}-\widehat R_{g+1}\) 的依赖。

把当前代码的两个未知量合并为
\begin{equation}
  \widehat{\mathbf x}_g
  =
  \begin{bmatrix}
  \widehat P_g\\
  \widehat R_g
  \end{bmatrix},
\end{equation}
则每个能量组满足
\begin{equation}
  \begin{bmatrix}
  1 & -B_g\\
  0 & C_g^R
  \end{bmatrix}
  \widehat{\mathbf x}_g
  +
  \begin{bmatrix}
  -C_g & C_g\\
  -H_g & H_g
  \end{bmatrix}
  \widehat{\mathbf x}_{g+1}
  =
  \begin{bmatrix}
  R_g^{(4)}\\
  N_g^0
  \end{bmatrix}.
\end{equation}
因此按 \(g=0,1,\ldots,N_g-1\) 排列，完整新值系统是 block 上双对角矩阵：
\begin{equation}
  \begin{bmatrix}
  A_0 & U_0 &        &        & 0\\
      & A_1 & U_1    &        &  \\
      &     & A_2    & \ddots &  \\
      &     &        & \ddots & U_{N_g-2}\\
  0   &     &        &        & A_{N_g-1}
  \end{bmatrix}
  \begin{bmatrix}
  \widehat{\mathbf x}_0\\
  \widehat{\mathbf x}_1\\
  \widehat{\mathbf x}_2\\
  \vdots\\
  \widehat{\mathbf x}_{N_g-1}
  \end{bmatrix}
  =
  \begin{bmatrix}
  \mathbf b_0\\
  \mathbf b_1\\
  \mathbf b_2\\
  \vdots\\
  \mathbf b_{N_g-1}
  \end{bmatrix},
\end{equation}
其中
\begin{equation}
  A_g=
  \begin{bmatrix}
  1 & -B_g\\
  0 & C_g^R
  \end{bmatrix},
  \qquad
  U_g=
  \begin{bmatrix}
  -C_g & C_g\\
  -H_g & H_g
  \end{bmatrix},
  \qquad
  \mathbf b_g=
  \begin{bmatrix}
  R_g^{(4)}\\
  N_g^0
  \end{bmatrix}.
\end{equation}
最高能边界取
\begin{equation}
  \widehat{\mathbf x}_{N_g}=0,
\end{equation}
所以这个 block 上双对角系统可以直接从 \(g=N_g-1\) 反向扫到 \(g=0\)，不需要
Thomas 算法，也不需要显式存储矩阵。代码中的
\begin{equation}
  \texttt{energyDgCnPrecomputeKernel}
  \quad\text{和}\quad
  \texttt{energyDgCnSolveFromPrecomputeKernel}
\end{equation}
正是先生成 \(A_g,U_g,\mathbf b_g\) 的等价系数，再做这个反向代入。

\section*{6. Transport Step: Eq. (18), Eq. (22)，当前代码未使用 CN}

固定能量组 \(g\) 和角度方向 \(m\)，transport 子步求解
\begin{equation}
  \partial_x \psi
  +
  \Omega_{m,y}\partial_y\psi
  +
  \Omega_{m,z}\partial_z\psi
  =
  0.
\end{equation}
这是论文 Eq. (18) 的 \(L_1\) 子问题；代码中 Eq. (22) 用 MUSCL 有限体积通量实现。
当前实现没有求解
\begin{equation}
  \psi^{n+1}
  =
  \psi^n-\frac{h}{2}
  \left[L_1(\psi^n)+L_1(\psi^{n+1})\right]
\end{equation}
这样的隐式/CN transport 系统，而是使用下面的显式 predictor--corrector。

MUSCL slope limiter 为
\begin{equation}
  \theta^y_{i,j}
  =
  \frac{\psi_{i,j}-\psi_{i-1,j}}
       {\psi_{i+1,j}-\psi_{i,j}},
  \qquad
  \eta(\theta)
  =
  \max\{0,\max[\min(2\theta,1),\min(\theta,2)]\},
\end{equation}
\begin{equation}
  \sigma^y_{i,j}
  =
  \eta(\theta^y_{i,j})(\psi_{i+1,j}-\psi_{i,j}),
  \qquad
  \sigma^z_{i,j}
  =
  \eta(\theta^z_{i,j})(\psi_{i,j+1}-\psi_{i,j}).
\end{equation}
界面重构为
\begin{equation}
  \psi^-_{i+1/2,j}=\psi_{i,j}+\frac12\sigma^y_{i,j},
  \qquad
  \psi^+_{i+1/2,j}=\psi_{i+1,j}-\frac12\sigma^y_{i+1,j},
\end{equation}
\begin{equation}
  \psi^-_{i,j+1/2}=\psi_{i,j}+\frac12\sigma^z_{i,j},
  \qquad
  \psi^+_{i,j+1/2}=\psi_{i,j+1}-\frac12\sigma^z_{i,j+1}.
\end{equation}
上风数值通量为
\begin{equation}
  H^y_{i+1/2,j}
  =
  \Omega_{m,y}
  \begin{cases}
    \psi^-_{i+1/2,j}, & \Omega_{m,y}\ge0,\\
    \psi^+_{i+1/2,j}, & \Omega_{m,y}<0,
  \end{cases}
\end{equation}
\begin{equation}
  H^z_{i,j+1/2}
  =
  \Omega_{m,z}
  \begin{cases}
    \psi^-_{i,j+1/2}, & \Omega_{m,z}\ge0,\\
    \psi^+_{i,j+1/2}, & \Omega_{m,z}<0.
  \end{cases}
\end{equation}
因此
\begin{equation}
  L_1(\psi_{i,j,g,m})
  =
  \frac{H^y_{i+1/2,j}-H^y_{i-1/2,j}}{\Delta y}
  +
  \frac{H^z_{i,j+1/2}-H^z_{i,j-1/2}}{\Delta z}.
\end{equation}
代码采用显式 predictor--corrector：
\begin{equation}
  \psi^\ast_{i,j,g,m}
  =
  \psi^n_{i,j,g,m}
  -
  h\,L_1(\psi^n_{i,j,g,m}),
\end{equation}
\begin{equation}
  \psi^{n+1}_{i,j,g,m}
  =
  \psi^n_{i,j,g,m}
  -
  \frac{h}{2}
  \left[
    L_1(\psi^n_{i,j,g,m})
    +
    L_1(\psi^\ast_{i,j,g,m})
  \right],
\end{equation}
其中 \(h=\Delta x/2\)。

\subsection*{6.1 Corrector Step 的代码离散式}

为了和代码变量一一对应，先把通量散度记为
\begin{equation}
  \mathcal L_{i,j,g,m}(\Psi)
  =
  \frac{H^y_{i+1/2,j,g,m}(\Psi)-H^y_{i-1/2,j,g,m}(\Psi)}{\Delta y}
  +
  \frac{H^z_{i,j+1/2,g,m}(\Psi)-H^z_{i,j-1/2,g,m}(\Psi)}{\Delta z}.
\end{equation}
其中 \(H^y,H^z\) 由上面的 MUSCL limiter 和 upwind 选择得到。
代码中的
\begin{equation}
  \texttt{d\_flux\_old}
  \quad\Longleftrightarrow\quad
  \mathcal L_{i,j,g,m}(\Psi^n),
\end{equation}
由
\begin{equation}
  \texttt{spatialFluxDivergencePlaneKernel(F, d\_flux\_old)}
\end{equation}
计算。

predictor 先给出临时值
\begin{equation}
  \Psi^\ast_{i,j,g,m}
  =
  \Psi^n_{i,j,g,m}
  -
  h\,\mathcal L_{i,j,g,m}(\Psi^n).
\end{equation}
代码对应
\begin{equation}
  \texttt{spatialPredictorKernel(F, d\_flux\_old, d\_F\_pred)}.
\end{equation}
然后用 \(\Psi^\ast\) 重新重构 MUSCL 界面值、重新计算上风通量，得到
\begin{equation}
  \texttt{d\_flux\_pred}
  \quad\Longleftrightarrow\quad
  \mathcal L_{i,j,g,m}(\Psi^\ast),
\end{equation}
代码对应
\begin{equation}
  \texttt{spatialFluxDivergencePlaneKernel(d\_F\_pred, d\_flux\_pred)}.
\end{equation}

corrector 真正写回的新值为
\begin{equation}
  \boxed{
  \Psi^{n+1}_{i,j,g,m}
  =
  \Psi^n_{i,j,g,m}
  -
  \frac{h}{2}
  \left[
    \mathcal L_{i,j,g,m}(\Psi^n)
    +
    \mathcal L_{i,j,g,m}(\Psi^\ast)
  \right].
  }
\end{equation}
这正是代码
\begin{equation}
  \texttt{spatialCorrectorKernel}
\end{equation}
中的
\begin{equation}
  \texttt{value = F[idx] - 0.5 * dt * (flux\_old[idx] + flux\_pred[idx])}.
\end{equation}
如果 \(\texttt{preserve\_sign=false}\)，代码还会做一次非负截断：
\begin{equation}
  \Psi^{n+1}_{i,j,g,m}
  \leftarrow
  \max\!\left(\Psi^{n+1}_{i,j,g,m},0\right).
\end{equation}
如果 \(\texttt{preserve\_sign=true}\)，则直接保留上式计算值。

\section*{7. Angle Diffusion Step: Eq. (19), Eq. (20), Eq. (21)，使用 CN}

固定 \((i,j,g)\)，角扩散子步求解论文 Eq. (19)：
\begin{equation}
  \partial_x\psi
  =
  D_g(\partial_{uu}\psi+\partial_{vv}\psi),
  \qquad
  D_g=\frac12\rho\,\Sigma_{\mathrm{tr},g}.
\end{equation}
代码使用零边界 sine basis，对应论文 Eq. (20), Eq. (21)。角网格内点为
\begin{equation}
  u_a=-\frac{L_u}{2}+(a+1)\Delta u,\qquad
  v_b=-\frac{L_v}{2}+(b+1)\Delta v.
\end{equation}
sine 模态为
\begin{equation}
  s^u_{\mu,a}=\sin\frac{\pi(\mu+1)(a+1)}{N_u+1},
  \qquad
  s^v_{\nu,b}=\sin\frac{\pi(\nu+1)(b+1)}{N_v+1}.
\end{equation}
离散 Laplacian 本征值为
\begin{equation}
  \lambda^u_\mu
  =
  \frac{2-2\cos\left(\frac{\pi(\mu+1)}{N_u+1}\right)}{\Delta u^2},
  \qquad
  \lambda^v_\nu
  =
  \frac{2-2\cos\left(\frac{\pi(\nu+1)}{N_v+1}\right)}{\Delta v^2}.
\end{equation}
这是当前实现中第二处使用 Crank--Nicolson 的地方。代码先用 sine/DST 将角扩散算子对角化，
然后对每个谱模态使用 CN 更新因子
\begin{equation}
  r_{g,\mu,\nu}
  =
  \frac{2/\Delta x-D_g(\lambda^u_\mu+\lambda^v_\nu)}
       {2/\Delta x+D_g(\lambda^u_\mu+\lambda^v_\nu)}.
\end{equation}
因此角扩散完整步为
\begin{equation}
  \widehat\psi_{g,\mu,\nu}
  =
  r_{g,\mu,\nu}\,\psi_{g,\mu,\nu},
\end{equation}
再做 inverse sine transform 回到角网格。

\subsection*{7.1 Forward / Inverse Sine Transform 的离散形式}

对固定的 \((y_i,z_j,E_g)\)，把角网格上的值记为
\begin{equation}
  \psi_{a,b}
  =
  \psi(y_i,z_j,u_a,v_b,E_g),
  \qquad
  a=0,\ldots,N_u-1,\quad b=0,\ldots,N_v-1.
\end{equation}
代码中的 sine basis 为
\begin{equation}
  S^u_{\mu,a}
  =
  \sin\frac{\pi(\mu+1)(a+1)}{N_u+1},
  \qquad
  S^v_{\nu,b}
  =
  \sin\frac{\pi(\nu+1)(b+1)}{N_v+1}.
\end{equation}

forward sine transform 在数学上可写为
\begin{equation}
  \widetilde\psi_{\mu,\nu}
  =
  \frac{4}{(N_u+1)(N_v+1)}
  r_{g,\mu,\nu}
  \sum_{b=0}^{N_v-1}
  \sum_{a=0}^{N_u-1}
  \psi_{a,b}
  S^u_{\mu,a}
  S^v_{\nu,b}.
\end{equation}
这里代码把归一化因子
\begin{equation}
  \frac{4}{(N_u+1)(N_v+1)}
\end{equation}
和 CN 更新因子 \(r_{g,\mu,\nu}\) 合并到了 forward 阶段的
\texttt{coeff\_factor} 中。

inverse sine transform 回到角网格为
\begin{equation}
  \psi^{n+1}_{a,b}
  =
  \sum_{\nu=0}^{N_v-1}
  \sum_{\mu=0}^{N_u-1}
  \widetilde\psi_{\mu,\nu}
  S^u_{\mu,a}
  S^v_{\nu,b}.
\end{equation}

因此，代码中小角度网格路径 \((N_u,N_v\le32)\) 的完整 angle step 是
\begin{equation}
\boxed{
  \psi^{n+1}_{a,b}
  =
  \sum_{\nu,\mu}
  \left[
    \frac{4}{(N_u+1)(N_v+1)}
    r_{g,\mu,\nu}
    \sum_{b',a'}
    \psi^n_{a',b'}
    S^u_{\mu,a'}
    S^v_{\nu,b'}
  \right]
  S^u_{\mu,a}
  S^v_{\nu,b}
}.
\end{equation}
代码对应关系为
\begin{equation}
\begin{array}{lll}
\texttt{sineForwardKernel} &:&
\psi_{a,b}\rightarrow \widetilde\psi_{\mu,\nu},\\
\texttt{buildDiffusionFactorKernel} &:&
r_{g,\mu,\nu},\\
\texttt{complexToRealKernel} &:&
\frac{4}{(N_u+1)(N_v+1)}r_{g,\mu,\nu},\\
\texttt{sineInverseKernel} &:&
\widetilde\psi_{\mu,\nu}\rightarrow \psi^{n+1}_{a,b}.
\end{array}
\end{equation}

\subsection*{7.2 大角度网格的推荐优化：Separable DST--GEMM}

当前大角度网格路径使用 odd extension 加 2D complex FFT：
\begin{equation}
  \text{pack odd extension}
  \rightarrow
  \text{2D C2C FFT}
  \rightarrow
  \text{multiply }r_{g,\mu,\nu}
  \rightarrow
  \text{2D inverse C2C FFT}
  \rightarrow
  \text{unpack}.
\end{equation}
该方法保持了 sine basis 的数学形式，但会把角网格从
\begin{equation}
  N_u\times N_v
\end{equation}
扩展为
\begin{equation}
  2(N_u+1)\times 2(N_v+1),
\end{equation}
并且使用 complex FFT，导致内存流量和 workspace 都较大。

由于角扩散算子是 separable 的，
\begin{equation}
  \partial_{uu}+\partial_{vv},
\end{equation}
二维 sine transform 可以拆成两个一维 sine transform。令
\begin{equation}
  \mathbf \Psi\in\mathbb R^{N_v\times N_u},
\end{equation}
其中行方向为 \(v\)，列方向为 \(u\)。定义 sine 矩阵
\begin{equation}
  (S_u)_{\mu,a}
  =
  \sin\frac{\pi(\mu+1)(a+1)}{N_u+1},
  \qquad
  (S_v)_{\nu,b}
  =
  \sin\frac{\pi(\nu+1)(b+1)}{N_v+1}.
\end{equation}
forward DST 可以写成矩阵形式：
\begin{equation}
  \widetilde{\mathbf \Psi}
  =
  \alpha
  S_v\,\mathbf \Psi\,S_u^T,
  \qquad
  \alpha=\frac{4}{(N_u+1)(N_v+1)}.
\end{equation}
CN 谱更新为。这里的 \(\mathrm{new}\) 是新深度步的状态标签，不表示幂次；为避免和幂次混淆，
下面统一把它写成下标。
\begin{equation}
  \widetilde{\mathbf \Psi}_{\mathrm{new},\nu,\mu}
  =
  r_{g,\mu,\nu}
  \widetilde{\mathbf \Psi}_{\nu,\mu}.
\end{equation}
inverse DST 为
\begin{equation}
  \mathbf \Psi_{\mathrm{new}}
  =
  S_v^T\,\widetilde{\mathbf \Psi}_{\mathrm{new}}\,S_u.
\end{equation}
合并起来就是
\begin{equation}
  \boxed{
  \mathbf \Psi_{\mathrm{new}}
  =
  S_v^T
  \left[
    \mathbf R_g
    \odot
    \left(
      \alpha S_v\,\mathbf \Psi\,S_u^T
    \right)
  \right]
  S_u
  },
\end{equation}
其中 \(\mathbf R_g=(r_{g,\mu,\nu})\)，\(\odot\) 表示逐元素乘。

实现上可以用 batched GEMM 或 fused shared-memory kernel：
\begin{equation}
\begin{aligned}
  \mathbf Y &= \mathbf \Psi S_u^T,\\
  \mathbf Z &= S_v \mathbf Y,\\
  \mathbf Z_{\nu,\mu} &\leftarrow \alpha r_{g,\mu,\nu}\mathbf Z_{\nu,\mu},\\
  \mathbf W &= S_v^T\mathbf Z,\\
  \mathbf \Psi_{\mathrm{new}} &= \mathbf W S_u.
\end{aligned}
\end{equation}
这里每个 batch 对应一个固定的 \((y_i,z_j,E_g)\) 平面。代码中的 batch 数为
\begin{equation}
  N_{\mathrm{batch}}
  =
  (N_y+1)(N_z+1)N_g.
\end{equation}

该优化的优点是：
\begin{itemize}
\item 不改变 Eq. (19)--(21) 的数学离散，只改变 sine transform 的实现方式。
\item 避免 odd extension，不再把 \(N_u\times N_v\) 扩展到 \(2(N_u+1)\times2(N_v+1)\)。
\item 避免 complex FFT，数据保持 real-valued。
\item 对 \(64\times64\)、\(128\times128\) 这类中大角网格，batched GEMM 可能比当前 odd-extension FFT 更快。
\item CN factor multiplication 可以和中间矩阵 \(\mathbf Z\) 的写回融合，减少一次全局内存读写。
\end{itemize}

建议的实现策略是保留现有路径作为 fallback：
\begin{equation}
\begin{aligned}
  &\text{small grid: direct DST}\\
  &\rightarrow\ \text{medium/large grid: separable DST--GEMM}\\
  &\rightarrow\ \text{very large grid: FFT or ADI fallback}.
\end{aligned}
\end{equation}
当前代码实现采用的分支为
\begin{equation}
\begin{array}{lll}
N_u,N_v\le 32 &:& \text{原 direct DST 路径},\\
33\le \max(N_u,N_v)\le 128 &:& \text{separable real DST 路径},\\
\max(N_u,N_v)>128 &:& \text{原 odd-extension FFT fallback}.
\end{array}
\end{equation}
当前 separable 路径先用四个 real-valued CUDA kernel 实现上述矩阵分解；
后续还可以进一步替换为 cuBLAS batched GEMM 或 shared-memory tiled kernel。

\subsection*{7.3 Separable DST 当前实现的详细步骤}

对每个固定的 \((y_i,z_j,E_g)\)，角平面是一个实矩阵
\begin{equation}
  \mathbf\Psi
  \in
  \mathbb R^{N_v\times N_u}.
\end{equation}
其中 \(N_u=\texttt{Nom}\)，\(N_v=\texttt{Nmu}\)。代码实际存储是 angle-major 的一维数组：
\begin{equation}
  \texttt{angle}=bN_u+a,
  \qquad
  a=0,\ldots,N_u-1,\quad b=0,\ldots,N_v-1.
\end{equation}
batch 编号对应
\begin{equation}
  \texttt{batch index}
  =
  g\,(N_y+1)(N_z+1)+\texttt{yz cell}.
\end{equation}

当前 separable real DST 路径按四个阶段执行。

\paragraph{\(Y,Z,W\) 与抽象 \(S,\Lambda\) 的对应关系}
前面写的
\begin{equation}
  L=S^{-1}\Lambda S
\end{equation}
是抽象的一维向量写法：把整个二维角平面 \(\Psi\) 拉直成向量
\(\boldsymbol\psi\)，然后用一个完整二维 DST 矩阵 \(S\) 做变换。
在这个抽象写法中，angle step 等价于
\begin{equation}
  \boldsymbol\psi_{\mathrm{new}}
  =
  S^{-1}R_gS\boldsymbol\psi,
\end{equation}
其中
\begin{equation}
  R_g
  =
  \left(I-\frac{\Delta x}{2}D_g\Lambda\right)^{-1}
  \left(I+\frac{\Delta x}{2}D_g\Lambda\right).
\end{equation}
因为 \(\Lambda\) 是对角矩阵，\(R_g\) 也是对角矩阵；其对角元就是代码中的
\begin{equation}
  r_{g,\mu,\nu}.
\end{equation}

代码没有显式构造完整二维矩阵 \(S\)。它利用二维 DST 的可分离性，把
\begin{equation}
  S\boldsymbol\psi
\end{equation}
写成两个方向的矩阵乘：
\begin{equation}
  S\boldsymbol\psi
  \quad\Longleftrightarrow\quad
  \alpha S_v\Psi S_u^T.
\end{equation}
其中 \(S_u\) 和 \(S_v\) 分别是一维 sine 矩阵，\(\alpha\) 是 DST 归一化因子。
同理，完整二维 inverse DST 不是单独的 \(S_v^T\)，而是
\begin{equation}
  S^{-1}(\cdot)
  \quad\Longleftrightarrow\quad
  S_v^T(\cdot)S_u.
\end{equation}
也就是说，\(S_v^T\) 只对应 \(v\) 方向的 inverse sine transform；右乘 \(S_u\)
对应 \(u\) 方向的 inverse sine transform。两步合起来才是完整的 \(S^{-1}\)。

于是中间矩阵 \(Y,Z,W\) 的含义是：
\begin{equation}
\begin{array}{lll}
\mathbf Y=\Psi S_u^T
&:& \text{完成 }u\text{ 方向 forward DST，是 }S\boldsymbol\psi\text{ 的第一半},\\[0.4em]
\mathbf Z=\mathbf R_g\odot(\alpha S_v\mathbf Y)
&:& \text{完成 }v\text{ 方向 forward DST，并乘 CN 因子},\\[0.4em]
\mathbf W=S_v^T\mathbf Z
&:& \text{完成 }v\text{ 方向 inverse DST，是 }S^{-1}\text{ 的第一半},\\[0.4em]
\Psi_{\mathrm{new}}=\mathbf W S_u
&:& \text{完成 }u\text{ 方向 inverse DST，得到新角分布}.
\end{array}
\end{equation}
因此
\begin{equation}
  \boxed{
  \Psi_{\mathrm{new}}
  =
  S_v^T
  \left[
    \mathbf R_g
    \odot
    \left(
      \alpha S_v\Psi S_u^T
    \right)
  \right]
  S_u
  }
\end{equation}
就是矩阵形式的
\begin{equation}
  \boxed{
  \boldsymbol\psi_{\mathrm{new}}
  =
  S^{-1}R_gS\boldsymbol\psi.
  }
\end{equation}
对应关系可以简写为
\begin{equation}
\begin{array}{lll}
S &\Longleftrightarrow& \alpha S_v(\cdot)S_u^T,\\[0.3em]
R_g &\Longleftrightarrow& \mathbf R_g\odot(\cdot),\\[0.3em]
S^{-1} &\Longleftrightarrow& S_v^T(\cdot)S_u.
\end{array}
\end{equation}
这里 \(\Lambda\) 没有显式出现，是因为它已经被吸收到 CN 更新因子
\begin{equation}
  r_{g,\mu,\nu}
  =
  \frac{2/\Delta x-D_g\lambda^+_{\mu,\nu}}
       {2/\Delta x+D_g\lambda^+_{\mu,\nu}}
\end{equation}
中了。

\paragraph{\(S_u\) 和 \(S_u^T\) 的含义}
定义
\begin{equation}
  (S_u)_{\mu,a}
  =
  \sin\frac{\pi(\mu+1)(a+1)}{N_u+1},
  \qquad
  \mu=0,\ldots,N_u-1,\quad a=0,\ldots,N_u-1.
\end{equation}
这里 \(a\) 是原始 \(u\) 网格点编号，\(\mu\) 是 \(u\) 方向 sine mode 编号。
因此
\begin{equation}
  S_u\in\mathbb R^{N_u\times N_u},
\end{equation}
其行对应 mode \(\mu\)，列对应 grid index \(a\)。于是
\begin{equation}
  (S_u^T)_{a,\mu}=(S_u)_{\mu,a}.
\end{equation}

因为角平面矩阵写成
\begin{equation}
  \mathbf\Psi\in\mathbb R^{N_v\times N_u},
\end{equation}
其中列方向是 \(u\)，所以右乘 \(S_u^T\)：
\begin{equation}
  \mathbf Y=\mathbf\Psi S_u^T
\end{equation}
就是对每一行 \(v_b\) 沿 \(u\) 方向做 forward sine transform。逐元素写为
\begin{equation}
  Y_{b,\mu}
  =
  \sum_{a=0}^{N_u-1}
  \Psi_{b,a}(S_u^T)_{a,\mu}
  =
  \sum_{a=0}^{N_u-1}
  \Psi_{b,a}
  \sin\frac{\pi(\mu+1)(a+1)}{N_u+1}.
\end{equation}
维度为
\begin{equation}
  (N_v\times N_u)(N_u\times N_u)=N_v\times N_u.
\end{equation}
同理，左乘 \(S_v\) 是沿 \(v\) 方向做 forward sine transform，左乘 \(S_v^T\)
是沿 \(v\) 方向做 inverse sine transform，右乘 \(S_u\) 是沿 \(u\) 方向做 inverse sine transform。

\paragraph{sine mode 编号的含义}
这里的 \(\mu\) 和 \(\nu\) 是 sine mode 编号，也就是正弦基函数的编号。
在 \(u\) 方向，原始网格点编号是
\begin{equation}
  a=0,\ldots,N_u-1,
\end{equation}
而 sine mode 编号是
\begin{equation}
  \mu=0,\ldots,N_u-1.
\end{equation}
第 \(\mu\) 个 mode 对应的基函数为
\begin{equation}
  \sin\frac{\pi(\mu+1)(a+1)}{N_u+1}.
\end{equation}
例如
\begin{equation}
\begin{array}{lll}
\mu=0 &:& \displaystyle \sin\frac{\pi(a+1)}{N_u+1},\\begin{equation}0.7em]
\mu=1 &:& \displaystyle \sin\frac{2\pi(a+1)}{N_u+1},\\begin{equation}0.7em]
\mu=2 &:& \displaystyle \sin\frac{3\pi(a+1)}{N_u+1}.
\end{array}
\end{equation}
低 mode 表示沿 \(u\) 方向变化较慢的正弦波，高 mode 表示变化较快的正弦波。
在 \(v\) 方向同理，\(\nu\) 是 \(v\) 方向的 sine mode 编号。

forward sine transform 的作用就是把角网格值
\begin{equation}
  \psi(u_a,v_b)
\end{equation}
转换成 mode 系数
\begin{equation}
  \widetilde\psi_{\mu,\nu}.
\end{equation}
使用 sine mode 的原因是：在零边界条件下，离散角 Laplacian
\begin{equation}
  \partial_{uu}+\partial_{vv}
\end{equation}
在 sine mode 空间中是对角的。因此 Eq. (19) 的 angle diffusion step 不需要解一个大的角空间矩阵，
而是对每个 mode \((\mu,\nu)\) 直接乘一个 CN 因子
\begin{equation}
  r_{g,\mu,\nu}.
\end{equation}

\paragraph{DST 和 CN 的区别}
DST 的全称是 Discrete Sine Transform，即离散正弦变换。它不是时间推进格式，
而是一种空间/角变量上的基变换：
\begin{equation}
  \psi_{a,b}
  \longrightarrow
  \widetilde\psi_{\mu,\nu}.
\end{equation}
在本问题中，DST 的作用是把零边界条件下的角 Laplacian 对角化。

CN 的全称是 Crank--Nicolson。它是 \(x\) 方向推进时对 angle diffusion 方程的深度离散格式。
如果直接在角网格空间写 CN，会得到矩阵方程
\begin{equation}
  \left(I-\frac{\Delta x}{2}D_g L\right)\psi^{n+1}
  =
  \left(I+\frac{\Delta x}{2}D_g L\right)\psi^n,
\end{equation}
其中 \(L\) 是角空间离散 Laplacian 矩阵。

使用 DST 后，相当于利用
\begin{equation}
  L = S^{-1}\Lambda S
\end{equation}
把这个矩阵系统对角化。于是 CN 更新不再需要求解大矩阵，而是在每个 sine mode 上做逐元素乘法：
\begin{equation}
  \widetilde\psi^{\,n+1}_{\mu,\nu}
  =
  r_{g,\mu,\nu}
  \widetilde\psi^{\,n}_{\mu,\nu}.
\end{equation}
因此二者分工是：
\begin{equation}
\begin{array}{lll}
\text{DST} &:& \text{变换工具，把角 Laplacian 对角化},\\
\text{CN} &:& \text{深度推进格式，给出每个 mode 的更新因子 }r_{g,\mu,\nu}.
\end{array}
\end{equation}
当前 separable 路径用矩阵乘法实现 DST；CN 本身在 mode 空间只是逐 mode 乘系数。

\paragraph{CN 矩阵方程的来源}
angle diffusion 子步的连续方程是
\begin{equation}
  \partial_x\psi
  =
  D_g(\partial_{uu}\psi+\partial_{vv}\psi).
\end{equation}
把角变量 \((u,v)\) 离散后，所有角网格点上的值排成向量
\begin{equation}
  \boldsymbol\psi
  =
  (\psi_{0,0},\psi_{1,0},\ldots,\psi_{N_u-1,N_v-1})^T.
\end{equation}
离散角 Laplacian 记为矩阵 \(L\)，即
\begin{equation}
  L\boldsymbol\psi
  \approx
  \partial_{uu}\psi+\partial_{vv}\psi.
\end{equation}
\(L\boldsymbol\psi\) 和原始的 \(\boldsymbol\psi\) 的关系是：\(\boldsymbol\psi\) 是当前角网格上的函数值，
而 \(L\boldsymbol\psi\) 是对这些函数值做二阶差分后得到的角向曲率/扩散项。
它不是另一个独立未知量，而是由当前 \(\boldsymbol\psi\) 线性计算出来的结果。
逐点写为
\begin{equation}
  (L\psi)_{a,b}
  =
  \frac{\psi_{a+1,b}-2\psi_{a,b}+\psi_{a-1,b}}{\Delta u^2}
  +
  \frac{\psi_{a,b+1}-2\psi_{a,b}+\psi_{a,b-1}}{\Delta v^2}.
\end{equation}
这里边界采用零边界条件；因此边界外侧的值按零处理。
直观上，如果 \(\psi\) 在角空间很平滑，\(L\psi\) 较小；如果 \(\psi\) 在角空间弯曲/变化很快，
\(L\psi\) 较大。angle diffusion step 正是用 \(D_gL\psi\) 来描述角分布在 \(u,v\) 方向的扩散。
于是得到关于深度变量 \(x\) 的 ODE 系统
\begin{equation}
  \frac{d\boldsymbol\psi}{dx}
  =
  D_gL\boldsymbol\psi.
\end{equation}

对这个 ODE 系统从 \(x^n\) 到 \(x^{n+1}=x^n+\Delta x\) 使用 Crank--Nicolson：
\begin{equation}
  \frac{\boldsymbol\psi^{n+1}-\boldsymbol\psi^n}{\Delta x}
  =
  \frac12D_gL\boldsymbol\psi^n
  +
  \frac12D_gL\boldsymbol\psi^{n+1}.
\end{equation}
两边乘以 \(\Delta x\) 并把含 \(\boldsymbol\psi^{n+1}\) 的项移到左边：
\begin{equation}
  \boldsymbol\psi^{n+1}
  -
  \frac{\Delta x}{2}D_gL\boldsymbol\psi^{n+1}
  =
  \boldsymbol\psi^n
  +
  \frac{\Delta x}{2}D_gL\boldsymbol\psi^n.
\end{equation}
因此得到矩阵方程
\begin{equation}
  \left(
    I-\frac{\Delta x}{2}D_gL
  \right)
  \boldsymbol\psi^{n+1}
  =
  \left(
    I+\frac{\Delta x}{2}D_gL
  \right)
  \boldsymbol\psi^n.
\end{equation}
这就是 angle step 中 CN 矩阵方程的来源。DST 的作用是把 \(L\) 对角化，从而避免直接求解这个大线性系统。

\paragraph{如何利用 \(L=S^{-1}\Lambda S\) 对角化 CN 系统}
从 CN 矩阵方程出发：
\begin{equation}
  \left(
    I-\frac{\Delta x}{2}D_gL
  \right)
  \boldsymbol\psi^{n+1}
  =
  \left(
    I+\frac{\Delta x}{2}D_gL
  \right)
  \boldsymbol\psi^n.
\end{equation}
设离散角 Laplacian 可以被 sine transform 对角化：
\begin{equation}
  L=S^{-1}\Lambda S,
\end{equation}
其中 \(S\) 表示二维 DST 变换，\(\Lambda\) 是对角矩阵：
\begin{equation}
  \Lambda
  =
  \mathrm{diag}(\lambda_{\mu,\nu}).
\end{equation}
定义 mode 空间变量
\begin{equation}
  \widetilde{\boldsymbol\psi}
  =
  S\boldsymbol\psi,
  \qquad
  \boldsymbol\psi
  =
  S^{-1}\widetilde{\boldsymbol\psi}.
\end{equation}
把
\begin{equation}
  \boldsymbol\psi^{n+1}=S^{-1}\widetilde{\boldsymbol\psi}^{\,n+1},
  \qquad
  \boldsymbol\psi^{n}=S^{-1}\widetilde{\boldsymbol\psi}^{\,n}
\end{equation}
代入 CN 方程：
\begin{equation}
  \left(
    I-\frac{\Delta x}{2}D_gS^{-1}\Lambda S
  \right)
  S^{-1}\widetilde{\boldsymbol\psi}^{\,n+1}
  =
  \left(
    I+\frac{\Delta x}{2}D_gS^{-1}\Lambda S
  \right)
  S^{-1}\widetilde{\boldsymbol\psi}^{\,n}.
\end{equation}
利用 \(SS^{-1}=I\)，并在方程两边左乘 \(S\)，得到
\begin{equation}
  \left(
    I-\frac{\Delta x}{2}D_g\Lambda
  \right)
  \widetilde{\boldsymbol\psi}^{\,n+1}
  =
  \left(
    I+\frac{\Delta x}{2}D_g\Lambda
  \right)
  \widetilde{\boldsymbol\psi}^{\,n}.
\end{equation}
这里需要注意：不是单独把 \(D_g\) 右侧的 \(S^{-1}\) 去掉，而是整个方程两边左乘 \(S\)
后成对相消。令
\begin{equation}
  a=\frac{\Delta x}{2}.
\end{equation}
以左边为例，先展开
\begin{equation}
  \left(I-aD_gS^{-1}\Lambda S\right)
  S^{-1}\widetilde{\boldsymbol\psi}^{\,n+1}
  =
  S^{-1}\widetilde{\boldsymbol\psi}^{\,n+1}
  -
  aD_gS^{-1}\Lambda SS^{-1}
  \widetilde{\boldsymbol\psi}^{\,n+1}.
\end{equation}
由于 \(SS^{-1}=I\)，上式变成
\begin{equation}
  S^{-1}\widetilde{\boldsymbol\psi}^{\,n+1}
  -
  aD_gS^{-1}\Lambda
  \widetilde{\boldsymbol\psi}^{\,n+1}.
\end{equation}
这时最前面的 \(S^{-1}\) 还在。接着对整个方程左乘 \(S\)，得到
\begin{equation}
  S
  \left[
    S^{-1}\widetilde{\boldsymbol\psi}^{\,n+1}
    -
    aD_gS^{-1}\Lambda
    \widetilde{\boldsymbol\psi}^{\,n+1}
  \right]
  =
  \widetilde{\boldsymbol\psi}^{\,n+1}
  -
  aD_g\Lambda
  \widetilde{\boldsymbol\psi}^{\,n+1}.
\end{equation}
这里 \(D_g\) 是标量，因此可与矩阵乘法交换；真正消掉的是
\begin{equation}
  SS^{-1}=I.
\end{equation}
右边同理，因此最后得到 mode 空间中的对角 CN 系统。
由于 \(\Lambda\) 是对角矩阵，这个系统已经完全解耦。对每个 mode \((\mu,\nu)\)，有
\begin{equation}
  \left(
    1-\frac{\Delta x}{2}D_g\lambda_{\mu,\nu}
  \right)
  \widetilde\psi^{\,n+1}_{\mu,\nu}
  =
  \left(
    1+\frac{\Delta x}{2}D_g\lambda_{\mu,\nu}
  \right)
  \widetilde\psi^{\,n}_{\mu,\nu}.
\end{equation}
因此
\begin{equation}
  \widetilde\psi^{\,n+1}_{\mu,\nu}
  =
  \frac{
    1+\frac{\Delta x}{2}D_g\lambda_{\mu,\nu}
  }{
    1-\frac{\Delta x}{2}D_g\lambda_{\mu,\nu}
  }
  \widetilde\psi^{\,n}_{\mu,\nu}.
\end{equation}
代码中使用的离散 Laplacian 特征值按正数
\begin{equation}
  \lambda^+_{\mu,\nu}
  =
  \lambda^u_\mu+\lambda^v_\nu
\end{equation}
存储，而扩散算子实际对应 \(-\lambda^+_{\mu,\nu}\)。因此代码里的 CN 因子写成
\begin{equation}
  r_{g,\mu,\nu}
  =
  \frac{
    2/\Delta x-D_g\lambda^+_{\mu,\nu}
  }{
    2/\Delta x+D_g\lambda^+_{\mu,\nu}
  }.
\end{equation}
这正是
\begin{equation}
  \texttt{buildDiffusionFactorKernel}
\end{equation}
中构造的谱空间更新因子。

\paragraph{第一步：沿 \(u\) 方向 forward DST}
对每个固定的 \(v_b\)，计算
\begin{equation}
  Y_{b,\mu}
  =
  \sum_{a=0}^{N_u-1}
  \Psi_{b,a} S^u_{\mu,a},
  \qquad
  \mu=0,\ldots,N_u-1.
\end{equation}
矩阵形式：
\begin{equation}
  \mathbf Y
  =
  \mathbf\Psi S_u^T,
  \qquad
  \mathbf Y\in\mathbb R^{N_v\times N_u}.
\end{equation}
代码对应：
\begin{equation}
  \texttt{separableForwardUKernel}.
\end{equation}

\paragraph{第二步：沿 \(v\) 方向 forward DST，并乘 CN 因子}
对每个固定的 \(u\)-mode \(\mu\)，计算
\begin{equation}
  Z_{\nu,\mu}
  =
  \alpha r_{g,\mu,\nu}
  \sum_{b=0}^{N_v-1}
  S^v_{\nu,b}Y_{b,\mu},
  \qquad
  \nu=0,\ldots,N_v-1.
\end{equation}
其中
\begin{equation}
  \alpha=\frac{4}{(N_u+1)(N_v+1)}
\end{equation}
是二维 DST-I 的归一化因子，
\begin{equation}
  r_{g,\mu,\nu}
  =
  \frac{2/\Delta x-D_g(\lambda^u_\mu+\lambda^v_\nu)}
       {2/\Delta x+D_g(\lambda^u_\mu+\lambda^v_\nu)}
\end{equation}
是 Eq. (19)--(21) 的 CN 谱更新因子。矩阵形式：
\begin{equation}
  \mathbf Z
  =
  \mathbf R_g
  \odot
  \left(
    \alpha S_v\mathbf Y
  \right).
\end{equation}
代码对应：
\begin{equation}
  \texttt{separableForwardVFactorKernel}.
\end{equation}
这里 \(\alpha r_{g,\mu,\nu}\) 预先存入 \texttt{d\_sine\_coeff}。

\paragraph{第三步：沿 \(v\) 方向 inverse DST}
对每个固定的 \(\mu\)，计算
\begin{equation}
  W_{b,\mu}
  =
  \sum_{\nu=0}^{N_v-1}
  S^v_{\nu,b} Z_{\nu,\mu}.
\end{equation}
矩阵形式：
\begin{equation}
  \mathbf W
  =
  S_v^T\mathbf Z,
  \qquad
  \mathbf W\in\mathbb R^{N_v\times N_u}.
\end{equation}
代码对应：
\begin{equation}
  \texttt{separableInverseVKernel}.
\end{equation}

\paragraph{第四步：沿 \(u\) 方向 inverse DST}
对每个固定的 \(v_b\)，计算
\begin{equation}
  \Psi_{\mathrm{new},b,a}
  =
  \sum_{\mu=0}^{N_u-1}
  W_{b,\mu}S^u_{\mu,a}.
\end{equation}
矩阵形式：
\begin{equation}
  \mathbf\Psi_{\mathrm{new}}
  =
  \mathbf W S_u.
\end{equation}
代码对应：
\begin{equation}
  \texttt{separableInverseUKernel}.
\end{equation}

综上，当前 separable 路径等价于
\begin{equation}
  \mathbf\Psi_{\mathrm{new}}
  =
  S_v^T
  \left[
    \mathbf R_g
    \odot
    \left(
      \alpha S_v \mathbf\Psi S_u^T
    \right)
  \right]
  S_u.
\end{equation}
它和原来的 sine/CN 离散数学上是一致的；区别只是没有通过 odd extension 和 complex FFT 实现。

\section*{8. Energy 和 Transport 是否解矩阵}

\subsection*{8.1 Energy Step 是否解矩阵}

energy step 对应 Eq. (15)，本质上来自一个沿能量组耦合的 DG/CN 离散系统。
如果把所有能量组同时写出来，它可以看成一个按能量方向耦合的线性代数系统。
但当前代码没有显式组装全局矩阵，也没有调用通用矩阵求解器。

当前实现利用能量方向的上风/三角结构，从高能组到低能组反向扫掠：
\begin{equation}
  g=N_g-1,N_g-2,\ldots,0.
\end{equation}
每个能量组只需要已知的高能邻居新值差
\begin{equation}
  \Delta_{g+1}
  =
  \widehat P_{g+1}-\widehat R_{g+1},
\end{equation}
然后局部求
\begin{equation}
  \widehat R_g
  =
  \frac{N_g^0+H_g\Delta_{g+1}}{C_g^R},
  \qquad
  \widehat P_g
  =
  B_g\widehat R_g+R_g^{(4)}+C_g\Delta_{g+1}.
\end{equation}
因此：
\begin{equation}
  \boxed{
  \text{Energy step 等价于求解一个能量方向的三角/递推系统，但代码没有显式组装矩阵。}
  }
\end{equation}
代码对应：
\begin{equation}
  \texttt{energyDgCnPrecomputeKernel}
  \quad+\quad
  \texttt{energyDgCnSolveFromPrecomputeKernel}.
\end{equation}

\subsection*{8.2 Transport Step 是否解矩阵}

transport step 当前不是矩阵求解。它对 Eq. (18), Eq. (22) 的 \(YZ\) 对流子步使用显式
MUSCL predictor--corrector：
\begin{equation}
  \psi^\ast
  =
  \psi^n-hL_1(\psi^n),
\end{equation}
\begin{equation}
  \psi^{n+1}
  =
  \psi^n
  -
  \frac{h}{2}
  \left[
    L_1(\psi^n)+L_1(\psi^\ast)
  \right].
\end{equation}
这里没有求解
\begin{equation}
  \psi^{n+1}
  =
  \psi^n
  -
  \frac{h}{2}
  \left[
    L_1(\psi^n)+L_1(\psi^{n+1})
  \right]
\end{equation}
这样的隐式/CN 空间耦合系统。因此：
\begin{equation}
  \boxed{
  \text{Transport step 当前不解矩阵，只做显式通量计算和显式更新。}
  }
\end{equation}
代码对应：
\begin{equation}
\begin{array}{c}
  \texttt{spatialFluxDivergencePlaneKernel}\\
  \downarrow\\
  \texttt{spatialPredictorKernel}\\
  \downarrow\\
  \texttt{spatialFluxDivergencePlaneKernel}\\
  \downarrow\\
  \texttt{spatialCorrectorKernel}.
\end{array}
\end{equation}

\section*{9. 当前实现中的重要差异}

\begin{itemize}
\item transport step 不使用 \(\sigma_c\)、\(S(E)\)、\(T(E)\)；这些只进入 Eq. (15) energy step 和 secondary source。
\item transport step 当前是显式 predictor--corrector，不是隐式迭代，也不是当前代码中的 CN 子步。
\item angle step 使用 Eq. (19)--(21) 的谱空间 CN 更新；传入的是完整步 \(\Delta x\)。
\item energy step 使用 Eq. (15) 的 DG/CN 递推；传入的是半步 \(\Delta x/2\)。
\item 因此，当前代码中实际使用 CN 的地方只有 \(\mathcal E\) energy step 和 \(\mathcal A\) angle diffusion step。
\end{itemize}

\end{document}
